% Generated from tex/chapters/ch04/06_構築の管理.tex for the SWEBOK structure tree
\subsubsection{構築計画}

構築計画では、構築方法、部品の作成順序、統合順序、統合戦略、品質管理プロセス、担当割り当て、必要なツールや環境を決める。構築方法の選択は、複雑さの低減、変化への対応、検証しやすさに影響する。

たとえば、最初から全体を一度に統合する方針を選べば、初期の計画は単純に見えるが、問題発見が遅れる可能性がある。小さな部品を段階的に統合する方針を選べば、統合のための足場やテスト環境が必要になるが、問題を早く見つけやすい。
